\documentclass{article}

\usepackage[utf8]{inputenc}
\usepackage[T1]{fontenc}
\usepackage{helvet}
\usepackage{geometry}
\usepackage{xcolor}
\usepackage{eso-pic}
\usepackage{fancyhdr}
\usepackage{parskip}
\usepackage{microtype}
\usepackage[english, ukrainian]{babel}
\usepackage{url}
\usepackage{amsmath}
\usepackage{amssymb}
\usepackage{amsthm}
\usepackage{longtable}
\usepackage{booktabs}
\usepackage{hyperref}
\usepackage{titlesec}

\definecolor{nato}{HTML}{293757}
\definecolor{footergray}{gray}{0.65}

\geometry{a4paper,left=3cm,right=3cm,top=3cm,bottom=3cm}
\renewcommand{\familydefault}{\sfdefault}
\setcounter{secnumdepth}{3}

\titleformat{\section}
  {\color{nato}\Huge\bfseries}{\thesection.}{1em}{}
\titlespacing*{\section}{0pt}{30pt}{12pt}

\titleformat{\subsection}
  {\color{nato}\LARGE\bfseries}{\thesubsection.}{1em}{}
\titlespacing*{\subsection}{0pt}{20pt}{8pt}

\titleformat{\subsubsection}
  {\color{nato}\Large\bfseries}{\thesubsubsection.}{1em}{}
\titlespacing*{\subsubsection}{0pt}{10pt}{6pt}

\fancyhf{}
\fancyhead[R]{\small\color{footergray} 31 березня 2026}
\fancyfoot[L]{\small\color{footergray} Технічні вимоги ЄСІКС. II. Господарська діяльність}

\fancyfoot[R]{\small\color{footergray}\thepage}

\newcommand\HUGE[1]{\fontsize{40}{40}\selectfont #1}

\renewcommand{\headrulewidth}{0pt}
\renewcommand{\footrulewidth}{0.4pt}
\renewcommand{\footrule}{\color{footergray}\hrule width \textwidth height 0.4pt}

\begin{document}

\pagestyle{fancy}

\begin{titlepage}
  \AddToShipoutPictureBG*{\AtPageLowerLeft{\color{nato}\rule{\paperwidth}{\paperheight}}}
  \color{white}\thispagestyle{empty}\vspace*{4cm}\centering
  {\Huge  \textbf{Технічні вимоги ЄСІКС} \par} \vspace{2cm}
  {\HUGE  \textbf{II. Господарська діяльність} \par} \vspace{2.5cm}
  {\LARGE \textbf{II-1. Бюджетне планування і контроль видатків (BUD)} \par}
  {\LARGE \textbf{II-2. Фінансовий і бухгалтерський облік (FIN)} \par}
  {\LARGE \textbf{II-3. Кадрова система (ACC)} \par}
  {\LARGE \textbf{II-4. Публічний документообіг і сервіс-деск (ITSM)} \par} \vspace{2.5cm}
  \vfill
  {\large 2026 \copyright\ Інформаційні Судові Системи \par}
\end{titlepage}

\newpage

\begin{center}
  {\Huge\color{nato}\bfseries Зміст\par}
\end{center}
\vspace{40pt}
\tableofcontents

\begin{abstract}
ТЕХНІЧНІ ВИМОГИ
НА СТВОРЕННЯ ЗАСОБУ ІНФОРМАТИЗАЦІЇ - КОМПОНЕНТІВ ГОСПОДАРСЬКОЇ ДІЯЛЬНОСТІ
ЄДИНОЇ СУДОВОЇ ІНФОРМАЦІЙНО-КОМУНІКАЦІЙНОЇ СИСТЕМИ.
\end{abstract}

\newpage
\section*{Вступна частина}

\subsection*{Умовні скорочення та визначення}

Терміни та скорочення, що використовуються в документі:

\begin{longtable}{p{4cm}p{10cm}}
\toprule
\textbf{Термін} & \textbf{Значення} \\
\midrule
Автентифікація & Електронна процедура, яка дає змогу підтвердити електронну ідентифікацію фізичної, юридичної особи, інформаційної або інформаційно-комунікаційної системи та/або походження та цілісність електронних даних \\
Авторизація & Частина процедури визначення та надання прав доступу до функцій та/або інформаційних ресурсів для роботи в інформаційній системі, після автентифікації \\
Валідація & Встановлення об’єктивних і задокументованих доказів того, що визначені вимоги до спеціально призначеного застосування можна послідовно виконати \\
API & Автоматизовані програмні інтеграції, відомі також як “прикладний інтерфейс програмування – application programming interface” \\
БД & База Даних — Організована колекція даних, яка зберігається та управляється з використанням комп'ютерної системи; забезпечує зберігання, організацію та доступ до даних з метою ефективного використання та обробки \\
БП & Бізнес Процес — Сукупність дій необхідних для здійснення людиною або системою, в результаті яких відбувається досягнення бажаного результату \\
КЕП & Кваліфікований електронний підпис (X.509) \\
РМК & Реєстраційно-моніторингова картка \\
АРМ & Автоматизоване робоче місце користувача \\
HSM & Hardware Secured Module \\
ЄСІКС & Єдина судова інформаційно-комунікаційна система \\
IAM & Identity and Access Management \\
SSO & Single Sign-On \\
BPMN & Business Process Model and Notation (ISO 19510) \\
ABAC & Attribute-Based Access Control (ISO 29146) \\
PEP & Policy Enforcement Point (ISO 29146) \\
PIP & Policy Information Point (ISO 29146) \\
PAP & Policy Administration Point (ISO 29146) \\
BUD & Бюджетне планування і контроль видатків \\
FIN & Фінансовий і бухгалтерський облік \\
ACC & Кадрова система \\
ITSM & Публічний документообіг і сервіс-деск \\
\bottomrule
\end{longtable}

\newpage
\subsection*{Загальні положення}

Технічні вимоги на основі «Концепції ЄСІКС» від 2024 року \cite{esiks} та ERP/1 Облік (Парус-подібної облікової системи).

\subsection*{Призначення та цілі впровадження}

З метою забезпечення ефективного управління господарською діяльністю судової системи в рамках ЄСІКС,
а також унеможливлення фрагментації підходів до бюджетування, фінансового обліку, кадрового управління
та сервісного обслуговування, є необхідним створення та впровадження централізованих бізнес-сервісів,
зокрема: бюджетного планування і контролю видатків (BUD), фінансового і бухгалтерського обліку (FIN),
кадрової системи (ACC) та публічного документообігу і сервіс-деску (ITSM).
Реалізація зазначених складових є передумовою побудови цілісної, масштабованої та керованої
архітектури господарської діяльності ЄСІКС, зниження витрат на розробку та супровід,
забезпечення єдиних стандартів обліку та користувацького досвіду.

\subsection*{Ревізія поточних систем}

При середній ціні на ліцензію в 40000 грк/рік мова може йти про екномію до 100 млн грн/рік при створенні
власного продукту, який покриває сутності і процеси існуючих облікових систем судової системи.

\begin{table}[htbp]
\centering
\caption{Структура розподілу постачальників ПЗ для групи модулів «II. Облік»}
\label{tab:oblik-suppliers}
\begin{tabular}{p{5.2cm} p{3.8cm} *{4}{c}}
\toprule
\textbf{Призначення} 
& \textbf{Назва системи} 
& \multicolumn{2}{c}{\textbf{Кількість робочих місць}} 
 \\
\cmidrule(lr){3-4}\cmidrule(lr){5-6}
& & \textbf{поточна} & \textbf{очікувана}  \\
\midrule
Управління персоналом 
& Кадри-Web 
& 1055 & 0 \\
Бухгалтерський облік 
& IS-PRO 
& 428 & 0 \\
Контроль за видатками 
& Казначейства 
& 110 & 0 \\
Бюджетне планування 
& АІС «ГРК-ВЕБ» 
& 187 & 0 \\
Податкова звітність 
& M.E.Doc 
& 267 & 0 \\
Інша фінансова звітність 
& Є-звітність 
& 160 & 0 \\
\bottomrule
\end{tabular}
\end{table}

\newpage
\section{Бюджетне планування \\ і контроль видатків}

\subsection{Вимоги до сутностей BUD}

\subsubsection{Фінансовий план}
Фінансовий план є основною реєстровою сутністю BUD, що агрегує бюджетні асигнування, кошториси та ліміти видатків на основі даних ERP/1.

\subsubsection{Коригування плану}
Коригування плану фіксує зміни до затвердженого фінансового плану з обов’язковим КЕП-підписанням та збереженням історії версій.

\subsubsection{Реєстр закупівель}
Реєстр закупівель містить плани державних закупівель, пов’язані з бюджетними позиціями та інтегровані з модулем CRM.

\subsubsection{Реєстр погоджень}
Реєстр погоджень зберігає статуси та історію бізнес-процесів узгодження бюджетних документів відповідно до BPMN.

\subsubsection{Реєстр кошторисів і асигнувань}
Кошториси та асигнування містять розписи за КЕКВ, паспорти бюджетних програм та бюджетні запити відповідно до АІС «ГРК-ВЕБ».

\subsubsection{Реєстр казначейських рахунків і зобов’язань}
Реєстр фіксує бюджетні зобов’язання, ліміти та відповідність платежів казначейським правилам контролю видатків.

\newpage
\subsection{Вимоги до процесів BUD}

\subsubsection{Планування бюджету}
Процес планування бюджету формує річний/квартальний фінансовий план з автоматичним контролем лімітів та прогнозуванням видатків.

\subsubsection{Коригування та контроль}
Процес коригування та контролю забезпечує внесення змін до плану з автоматичною валідацією на перевищення асигнувань та генерацією звітів.

\subsubsection{Погодження видатків}
Процес погодження видатків реалізує багатоступеневий BPMN-роутінг документів з ролями ABAC та фіксацією всіх переходів у РМК.

\subsubsection{Формування бюджетних запитів і декларацій}
Процес формує бюджетні запити, декларації та паспорти програм з автоматичним контролем лімітів і експортом до АІС «ГРК-ВЕБ».

\subsubsection{Казначейський контроль платежів}
Процес забезпечує попередній та поточний контроль відповідності платежів асигнуванням, КЕКВ і казначейським рахункам з фіксацією в РМК.

\subsubsection{Інтеграція з Державним казначейством}
Процес здійснює автоматичний обмін розписами, зобов’язаннями та звітами з системою Казначейства.

\subsection{Вимоги до сховища та серіалізації BUD}
Сховище BUD використовує реляційне SQL для облікових сутностей та KV-простір для оперативних балансів; серіалізація здійснюється в JSON з JSON Schema валідацією.

\newpage
\section{Фінансовий і бухгалтерський облік}

\subsection{Вимоги до сутностей FIN}

\subsubsection{Головна книга}
Головна книга є центральним регістром бухгалтерських рахунків і проводок відповідно до плану рахунків ERP/1.

\subsubsection{Первинні документи}
Первинні документи містять реквізити господарських операцій з обов’язковим КЕП та інтеграцією з документообігом.

\subsubsection{Журнал проводок}
Журнал проводок фіксує всі бухгалтерські записи з автоматичною генерацією балансу та оборотних відомостей.

\subsubsection{Звіти та аналітика}
Звіти та аналітика забезпечують формування балансу, звіту про фінансові результати та аналітичних дашбордів.

\subsubsection{Картотека основних засобів і НМА}
Картотека ОЗ/НМА містить повний облік руху, амортизації, переоцінки, інвентаризації та податкового обліку відповідно до IS-PRO.

\subsubsection{Податкові регістри та ПДВ}
Податкові регістри забезпечують облік ПДВ, податкових накладних та звітності за формами ДПС.

\subsubsection{Форми фінансової звітності}
Форми фінансової звітності відповідають П(С)БО та Є-звітності (баланс, звіт про фінансові результати, консолідовані форми).

\newpage
\subsection{Вимоги до процесів FIN}

\subsubsection{Проведення операцій}
Процес проведення операцій виконує подвійний запис у головній книзі з автоматичною валідацією та закриттям періоду.

\subsubsection{Закриття періоду}
Процес закриття періоду формує звітність за місяць/квартал/рік з блокуванням змін після підписання КЕП.

\subsubsection{Автоматична звітність}
Процес автоматичної звітності генерує регламентовані форми для органів контролю та архівує їх у модулі CRM.

\subsubsection{Подвійний запис за П(С)БО}
Процес виконує подвійний запис з автоматичною валідацією за національними стандартами бухгалтерського обліку.

\subsubsection{Розрахунок амортизації та переоцінка}
Процес нараховує бухгалтерську, податкову та управлінську амортизацію ОЗ/НМА з веденням архіву змін.

\subsubsection{Формування та подання Є-звітності}
Процес генерує та підписує КЕП звіти для Казначейства та Мінфіну з автоматичним завантаженням у Є-звітність.

\subsection{Вимоги до сховища та серіалізації FIN}
Сховище FIN використовує SQL для транзакцій та S3 для сканованих первинних документів; серіалізація здійснюється в JSON з обов’язковою аудитом подій.

\addtocontents{toc}{\protect\newpage}
\newpage
\section{Кадрова система}

\subsection{Вимоги до сутностей ACC}

\subsubsection{Кадровий реєстр}
Кадровий реєстр містить дані про співробітників, посади, табельні номери та історії призначень.

\subsubsection{Розрахунок заробітної плати}
Розрахунок заробітної плати агрегує табелі, нарахування, утримання та податкові зобов’язання.

\subsubsection{Табелі обліку часу}
Табелі обліку часу фіксують відпрацьований час, відпустки та лікарняні з інтеграцією в процеси BUD.

\subsubsection{Звіти з кадрів}
Звіти з кадрів забезпечують формування відомостей про фонд оплати праці та кадрову статистику.

\subsubsection{Суддівські досьє та кар’єра}
Суддівські досьє містять дані про кар’єру, стаж, адміністративні посади, відсторонення та взаємодію з ВККСУ/ВРП.

\subsubsection{Штатний розпис і накази}
Штатний розпис і кадрові накази фіксують прийом/переміщення/звільнення працівників апарату та суддів.

\subsubsection{Табелі та розрахунок ЗП з нарахуваннями}
Табелі обліку часу інтегруються з фондом оплати праці та бюджетними лімітами BUD.

\newpage
\subsection{Вимоги до процесів ACC}

\subsubsection{Прийом та переміщення}
Процес прийому та переміщення реалізує кадрові накази з BPMN-узгодженням та автоматичним оновленням ролей ABAC.

\subsubsection{Розрахунок та виплата}
Процес розрахунку та виплати виконує нарахування заробітної плати з контролем бюджетних лімітів BUD.

\subsubsection{Формування звітності}
Процес формування звітності генерує податкові та статистичні звіти з КЕП-підписанням.

\subsubsection{Кадрове забезпечення суддівської кар’єри}
Процес оформлює зарахування/відрахування суддів, обчислення стажу та підготовку матеріалів для ВККСУ.

\subsubsection{Автоматизований кадровий документообіг}
Процес формує накази, графіки відпусток і дистанційної роботи з КЕП-підписанням та інтеграцією з Кадри-Web.

\subsubsection{Інтеграція з автоматизованим розподілом справ}
Процес передає актуальні кадрові дані до модуля авторозподілу справ ЄСІКС.

\subsection{Вимоги до сховища та серіалізації ACC}
Сховище ACC використовує SQL для кадрових даних та KV для оперативних розрахунків; серіалізація здійснюється в JSON з шифруванням персональних даних.

\newpage
\section{Публічний документообіг і сервіс-деск (ITSM)}

\subsection{Вимоги до сутностей ITSM}
Згідно з політиками проектування, у складі підсистеми публічного документообігу та сервіс-деску (ITSM) повинні підтримуватися наступні реєстрові сутності:
\begin{itemize}
    \item \textbf{Сервіс (Service)} — опис ІТ-послуги з каталогу послуг.
    \item \textbf{Звернення (Request / REQ)} — первинне звернення користувача через єдину точку контакту (SPOC).
    \item \textbf{Інцидент (Incident / INC)} — процесна сутність, що реєструє незаплановане припинення або деградацію ІТ-послуги.
    \item \textbf{Запит на обслуговування (Request for Service / RFS)} — звернення на надання інформації, доступу чи стандартних робіт.
    \item \textbf{Запит на зміну (Request for Change / RFC / CHG)} — процесна сутність, яка описує плановану або аварійну зміну в ІТ-середовищі.
    \item \textbf{Конфігураційний елемент (Configuration Item / CI)} — будь-який компонент ІТ-інфраструктури, що перебуває під управлінням конфігураціями.
    \item \textbf{Угода про рівень послуг (Service Level Agreement / SLA)} — параметри якості послуги, терміни реакції та вирішення.
\end{itemize}

\subsection{Вимоги до процесів та їх кроків}
Підсистема ITSM повинна реалізовувати наступні регламентні процеси та відповідні процедури:

\subsubsection{Управління сервісним каталогом (SCM)}
Процес забезпечує актуальність каталогу ІТ-послуг та включає такі кроки:
\begin{enumerate}
    \item Створення попереднього сервісу (чернетки).
    \item Редагування попереднього сервісу.
    \item Публікація сервісу (активація в каталозі).
    \item Деактивація сервісу.
    \item Архівування сервісу.
\end{enumerate}

\subsubsection{Управління інцидентами (INC)}
Процес спрямований на найшвидше відновлення працездатності послуг та включає такі кроки:
\begin{enumerate}
    \item Ініціація та реєстрація звернення користувача (REQ-P01).
    \item Первинна обробка звернення (тріаж) та ініціація інциденту (REQ-P02).
    \item Прийняття інциденту до опрацювання (INC-P01).
    \item Опрацювання інциденту (діагностика, вирішення або ескалація) (INC-P02).
    \item Закриття звернення після підтвердження вирішення (REQ-P03).
    \item Оперативний контроль, звітність та моніторинг показників (INC-P03).
    \item Оцінка та вдосконалення процесу (INC-P04).
\end{enumerate}

\subsubsection{Управління змінами (CHG)}
Процес мінімізує ризики при внесенні змін до ІТ-середовища та включає такі кроки:
\begin{enumerate}
    \item Ініціація та реєстрація потреби у зміні (REQ-P01).
    \item Первинна обробка звернення та ініціація запиту на зміну (REQ-P02).
    \item Прийняття зміни до опрацювання (CHG-P01).
    \item Аналіз та оцінка ризиків зміни (CHG-P02).
    \item Підготовка плану реалізації, тестування та плану повернення/відкату (CHG-P03).
    \item Реалізація зміни та верифікація результатів (CHG-P04).
    \item Закриття звернення (REQ-P03).
    \item Оперативний контроль і звітність процесу (CHG-P05).
    \item Оцінка та вдосконалення процесу (CHG-P06).
\end{enumerate}

\subsubsection{Управління конфігураціями (SACM)}
Процес забезпечує точність даних про ІТ-інфраструктуру та включає такі кроки:
\begin{enumerate}
    \item Ініціація та планування конфігураційного обліку (CFG-P01).
    \item Ідентифікація та реєстрація конфігураційного елемента в CMDB (CFG-P02).
    \item Контроль, моніторинг та облік стану конфігураційного елемента (CFG-P03).
    \item Верифікація, аудит та звірка фактичного стану з CMDB (CFG-P04).
\end{enumerate}

\subsubsection{Управління рівнем обслуговування (SLM)}
Процес визначає та контролює дотримання SLA для послуг та включає такі кроки:
\begin{enumerate}
    \item Визначення та фіксація параметрів SLA для послуг (SLM-P01).
    \item Моніторинг дотримання встановлених параметрів SLA (SLM-P02).
    \item Підтримка актуальності параметрів SLA (SLM-P03).
    \item Оперативний контроль і звітність процесу (SLM-P04).
    \item Оцінка та вдосконалення процесу (SLM-P05).
\end{enumerate}

\subsection{Вимоги до прав та привілеїв (PEP/ABAC)}
Доступ до операцій підсистеми ITSM керується на основі атрибутів суб'єктів та об'єктів (ABAC). Визначаються наступні обов'язкові точки перевірки правил (PEP) та відповідні права:
\begin{itemize}
    \item \textbf{VIEW\_ITSM} — право на перегляд реєстрів звернень, сервісів та інцидентів.
    \item \textbf{CARD\_ITSM} — право на перегляд детальної картки сутності.
    \item \textbf{CREATE\_ITSM} — право на створення звернень, запитів або конфігураційних елементів.
    \item \textbf{EDIT\_ITSM} — право на редагування карток та параметрів.
    \item \textbf{SCM\_PUBLISH} — привілей публікації та деактивації послуг у каталозі.
    \item \textbf{INC\_ASSIGN} — право на призначення та перенаправлення інциденту.
    \item \textbf{INC\_RESOLVE} — право на фіксацію рішення по інциденту.
    \item \textbf{CHG\_ANALYZE} — право на проведення аналізу ризиків та планування змін.
    \item \textbf{CHG\_EXECUTE} — право на реалізацію та закриття змін.
    \item \textbf{CFG\_AUDIT} — право на проведення верифікації та аудиту конфігураційних даних.
    \item \textbf{SLM\_DEFINE} — право на створення та зміну параметрів SLA.
\end{itemize}

\subsection{Вимоги до сховища та серіалізації ITSM}
З метою забезпечення прогнозованого навантаження та високої надійності:
\begin{itemize}
    \item Для зберігання структурованих реєстрових та реляційних даних (Сервіси, Зв'язки конфігурацій, SLA параметри) повинна використовуватися реляційна база даних (SQL).
    \item Для зберігання історії процесів, журналів переходів станів та текстових логів повідомлень (звернень) повинні використовуватися низьколатентні сховища з єдиним простором ключів (KV) на базі LSM-дерев та кешуванням у пам'яті.
    \item Взаємодія між компонентами та зовнішня серіалізація повинна здійснюватися у форматі JSON із обов'язковою автоматичною перевіркою схем за допомогою JSON Schema.
\end{itemize}

\newpage
\section*{Висновок}

Дотримання принципу «політики → вимоги → технічне завдання» для модулів господарської діяльності (BUD, FIN, ACC, ITSM) на основі ERP/1 Облік забезпечує повну інтеграцію з інфраструктурними сервісами ЄСІКС, технологічну незалежність та відповідність державним стандартам ДСТУ.

\begin{thebibliography}{9}

\bibitem{ecoz:rehab} Міністерство охорони здоров’я України. \newblock Технічні вимоги Реабілітація 2.0. \newblock Версія 08.08.2024.
\bibitem{dstu:3973} ДСТУ 3973-2000. \newblock Система розроблення та постачання продукції на виробництво. Правила виконання науково-дослідних робіт. Загальні положення.
\bibitem{dstu:3974} ДСТУ 3974-2000. \newblock Система розроблення та постачання продукції на виробництво. Правила виконання науково-конструкторських робіт. Загальні положення.
\bibitem{dstu:42010} ДСТУ 42010:2018. \newblock Інженерія систем і програмних засобів. Опис архітектури.
\bibitem{law:info-protection} Закон України «Про захист інформації в інформаційно-телекомунікаційних системах».
\bibitem{law:personal-data} Закон України «Про захист персональних даних».
\bibitem{law:cybersecurity} Закон України «Про основні засади забезпечення кібербезпеки України».
\bibitem{owasp:top10} OWASP Foundation. \newblock OWASP Top 10: The Ten Most Critical Web Application Security Risks. \newblock \url{https://owasp.org/www-project-top-ten/}.
\bibitem{nist:sp800-57} NIST Special Publication 800-57 Part 1 Revision 5. \newblock Recommendation for Key Management.
\bibitem{iso:42010} ISO/IEC/IEEE 42010:2022. \newblock Systems and software engineering — Architecture description.
\bibitem{zachman} Zachman, J.A. \newblock A Framework for Information Systems Architecture. \newblock IBM Systems Journal, 1987.
\bibitem{esiks} Концепція ЄСІКС. \url{https://court.gov.ua/storage/portal/dsa/news/Kontseptsia%20ESIKS.pdf}. 2024

\end{thebibliography}

\end{document}
